\section{Thinness does not imply splitting without chamber connectivity}
The last observation in Section 8 of the source is stated for thin,
flag-transitive geometries. Its definition of a geometry requires every
flag to extend to a chamber, but does not require the chamber graph to
be connected. The distinction is necessary.

We use one completely specified finite bipartite graph $B$. Its vertex
set is $\{1,\ldots,896\}$; the complete adjacency lists, together with
explicit permutation generators, are supplied in the mandatory file
\texttt{graph.json}. No graph catalogue or search is a premise.
The complete certificate program checks the following finite facts:
\begin{enumerate}
\item $B$ is simple, connected and cubic, with parts of size $448$ and
$1344$ edges.
\item The displayed permutations preserve adjacency and generate a
subgroup $H$ of order $5376$.
\item Starting with vertex $1$ and its three neighbors in four separate
singleton colors, and all other vertices in a fifth color, repeated
neighbor-color refinement has numbers of cells
\[
5,8,11,14,17,20,26,44,135,398,692,842,890,896.
\]
\item The subgroup $H^+$ preserving the two parts has order $2688$ and
one orbit on the $1344$ unordered edges.
\item The $2688$ elements exchanging the parts have orders
\[
\begin{array}{c|rrr}
\text{order}&4&12&28\\\hline
\text{number}&1024&896&768.
\end{array}
\]
\end{enumerate}
These are exhaustive finite calculations on the supplied graph and
permutations, rather than tests of a sample or of a proposed graph family.
Here is why they prove fullness without relying on a graph-automorphism
package. At each refinement step a vertex receives its previous color
and the multiset of colors of its neighbors. Every automorphism fixing
the four initially distinguished vertices fixes every resulting color.
The final partition is discrete, so the pointwise stabilizer of that
star is trivial. A vertex stabilizer therefore embeds in the symmetric
group on its three neighbors and has order at most $6$. Consequently
\[
 |\Aut(B)|\le 896\cdot6=5376.
\]
The explicitly generated subgroup attains this bound, so $H=\Aut(B)$.
The final two calculations concern the full automorphism group.

\begin{theorem}[A finite thin flag-transitive nonsplit example]
\label{A126.ThinCounterexample}
There is a finite rank-three thin, flag-transitive geometry $\Delta$
with $4480$ objects and $10752$ chambers whose full correlation group
does not split over its type-preserving subgroup. More precisely, for
$N=C_2^{2240}$ and the full groups just certified,
\[
 C(\Delta)=N\rtimes H,\qquad K(\Delta)=N\rtimes H^+,
 \qquad C(\Delta)/K(\Delta)\cong C_2.
\]
The chamber graph has exactly $1344$ connected components.
\end{theorem}
\begin{proof}
Apply the thin-doubling proposition of the common paper to $B$.
Its type sizes are $896,896,2688$, so the edge type is recognized.
There are two objects over each of the $896+1344=2240$ underlying
vertices and edges, and eight chambers over each edge. The proposition
therefore gives all the stated sizes and the exact full group pair.
A complement to $K(\Delta)$ would contain an involution in the
exchanging coset. Its projection to $H$ would be an exchanging
involution, contradicted by the exhaustive order table. This proves
nonsplitting. The connected-case splitting theorem of the common paper
is unaffected: each of the $1344$ edge cubes is one chamber component.
\end{proof}

This refutes the unqualified geometric splitting observation, under the
source's stated definition. It is not a classification of groups admitting
thin exact representations, nor a solution of the broader geometric
representation questions in Section 8. The graph certificate is a required
finite premise of this example. The uniform doubling and connected-case
arguments are written proofs; no finite sample substitutes for them.
